\section{Problem 6: Dominant Element}

\subsection{Problem Overview and Objectives}
Given an integer array of size $n$, determine whether there exists a \textbf{dominant element}—an element that occurs more than half of the time ($> n / 2$).

The objectives are:
\begin{enumerate}
    \item Develop **Solution 1**: Repeated Counting / Comparison (Brute Force).
    \item Develop **Solution 2**: Boyer-Moore Majority Vote Algorithm (Optimal linear time and constant memory).
    \item Provide mathematical proof explaining how the dominant element can be identified in linear time $O(n)$ \textbf{without using additional storage} ($O(1)$ space).
\end{enumerate}

\subsection{Specifications}
\begin{itemize}
    \item \textbf{Input}: An array of $n$ integers \texttt{numbers}.
    \item \textbf{Output}: \texttt{(has\_dominant, candidate\_value, occurrence\_count)}.
    \item \textbf{Constraints}: Solution 2 must allocate zero extra arrays, sets, or hash tables.
\end{itemize}

\subsection{Identifying Dominant Element Without Additional Storage}
\begin{quote}
\textbf{Question}: Can the dominant element be identified without using additional storage? \\
\textbf{Answer}: \textbf{YES.}
\end{quote}

\textbf{Mathematical Proof / Pairwise Cancellation Principle}:
\begin{enumerate}
    \item \textbf{Pairwise Cancellation Concept}: Suppose we repeatedly pair up and eliminate two \emph{different} elements from the array. Each elimination removes at most one instance of the dominant element and one instance of a non-dominant element.
    \item \textbf{Guaranteed Survival}: If an element $X$ appears $c > \lfloor n / 2 \rfloor$ times, the aggregate count of all other non-$X$ elements is $n - c < n/2$. Even if every non-$X$ element is paired with an instance of $X$ and eliminated, there will remain at least:
    \[
    c - (n - c) = 2c - n \ge 1
    \]
    un-eliminated instance of $X$.
    \item \textbf{Phase 1 (Candidate Search)}: Boyer-Moore tracks a single \texttt{candidate} variable and a scalar \texttt{vote\_counter}. When \texttt{vote\_counter} drops to 0, a new candidate is selected. By the cancellation principle, the true dominant element (if it exists) is guaranteed to be the surviving candidate at the end of Phase 1.
    \item \textbf{Phase 2 (Verification Pass)}: A second linear $O(n)$ pass counts exact occurrences of \texttt{candidate} to confirm if count $> \lfloor n / 2 \rfloor$.
\end{enumerate}

\subsection{Formal Algorithm Design}

\begin{algorithm}[H]
\caption{Solution 1: Brute Force Repeated Counting ($O(n^2)$ Time)}
\begin{algorithmic}[1]
\Function{FindDominantRepeatedCounting}{$numbers$}
    \State $n \gets \text{Length}(numbers)$
    \If{$n == 0$} \Return $(\text{False}, \text{None}, 0)$ \EndIf
    \State $threshold \gets \lfloor n / 2 \rfloor$
    \For{$i \gets 0$ \textbf{to} $n - 1$}
        \State $candidate \gets numbers[i]$, $count \gets 0$
        \For{$j \gets 0$ \textbf{to} $n - 1$}
            \If{$numbers[j] == candidate$}
                \State $count \gets count + 1$
            \EndIf
        \EndFor
        \If{$count > threshold$}
            \State \Return $(\text{True}, candidate, count)$
        \EndIf
    \EndFor
    \State \Return $(\text{False}, \text{None}, 0)$
\EndFunction
\end{algorithmic}
\end{algorithm}

\begin{algorithm}[H]
\caption{Solution 2: Boyer-Moore Majority Vote Algorithm ($O(n)$ Time, $O(1)$ Space)}
\begin{algorithmic}[1]
\Function{FindDominantBoyerMoore}{$numbers$}
    \State $n \gets \text{Length}(numbers)$
    \If{$n == 0$} \Return $(\text{False}, \text{None}, 0)$ \EndIf
    \State $candidate \gets \text{None}$, $votes \gets 0$ \Comment{Phase 1: Candidate Search}
    \For{\textbf{each} $num$ \textbf{in} $numbers$}
        \If{$votes == 0$}
            \State $candidate \gets num$, $votes \gets 1$
        \ElsIf{$num == candidate$}
            \State $votes \gets votes + 1$
        \Else
            \State $votes \gets votes - 1$
        \EndIf
    \EndFor
    \State $actual\_count \gets 0$ \Comment{Phase 2: Candidate Verification}
    \For{\textbf{each} $num$ \textbf{in} $numbers$}
        \If{$num == candidate$}
            \State $actual\_count \gets actual\_count + 1$
        \EndIf
    \EndFor
    \If{$actual\_count > \lfloor n / 2 \rfloor$}
        \State \Return $(\text{True}, candidate, actual\_count)$
    \Else
        \State \Return $(\text{False}, \text{None}, 0)$
    \EndIf
\EndFunction
\end{algorithmic}
\end{algorithm}

\subsection{Complexity Analysis and Comparison}

\subsubsection{Solution 1 (Brute Force Repeated Counting)}
\begin{itemize}
    \item \textbf{Best-Case Time Complexity}: $O(n)$.
    Occurs when the first element $numbers[0]$ is dominant and its frequency scan immediately exceeds $n/2$.
    \item \textbf{Worst-Case Time Complexity}: $O(n^2)$.
    Occurs when no dominant element exists or the dominant element is at the end. Outer loop runs $n$ times, inner loop runs $n$ times $\Rightarrow n \times n = n^2$ comparisons.
    \item \textbf{Auxiliary Space Complexity}: $O(1)$.
\end{itemize}

\subsubsection{Solution 2 (Boyer-Moore Majority Vote Algorithm)}
\begin{itemize}
    \item \textbf{Best, Average, \& Worst-Case Time Complexity}: $O(n)$.
    Phase 1 performs 1 pass over array ($n$ steps). Phase 2 performs 1 pass over array ($n$ steps). Total operations $= 2n = O(n)$ across all input distributions.
    \item \textbf{Auxiliary Space Complexity}: $O(1)$.
    Uses only two scalar variables (\texttt{candidate} and \texttt{votes}).
\end{itemize}

\begin{table}[h]
\centering
\caption{Comparison Matrix for Dominant Element Strategies}
\begin{tabular}{lcccc}
\toprule
\textbf{Approach} & \textbf{Best Time} & \textbf{Average Time} & \textbf{Worst Time} & \textbf{Auxiliary Space} \\
\midrule
\textbf{1. Repeated Counting (Brute Force)} & $O(n)$ & $O(n^2)$ & $O(n^2)$ & $O(1)$ \\
\textbf{2. Boyer-Moore Majority Vote} & $O(n)$ & $O(n)$ & $O(n)$ & $O(1)$ \\
\bottomrule
\end{tabular}
\end{table}

\textbf{Conclusion}: Solution 2 (Boyer-Moore) is strictly superior as it achieves optimal linear $O(n)$ time while strictly satisfying the zero additional storage constraint ($O(1)$ auxiliary space).
